\subsection{Entropy for a prescribed ordinary type}

The smaller-background condition lets the preceding entropy bounds survive prescribed singularities. The comparison uses a bounded subsolution in the part of the background left over by that condition.

\begin{corollary}\label{cor:prescribed-dual-entropy}
Under the standing hypotheses \eqref{eq:normalization-reference}, normalize $\mathrm{d}V(X)=1$. For every $q>1$ there is a finite real $a_q$, and there is a finite real $A$, such that
\[
 a_q+\frac1q\operatorname{Ent}_{\mathrm{d}V}(\nu)
 \leq H_\varphi(\nu)
 \leq A+n\operatorname{Ent}_{\mathrm{d}V}(\nu)
\]
for every Borel probability $\nu$, including infinite values. Thus finite dual value and finite entropy have the same domain, and uniformly bounded dual values are equivalent to uniformly bounded entropies. For fixed $\theta$ and $\varepsilon$ in \eqref{eq:normalization-reference}, the constants can be chosen independently of the model reference $\varphi$.
\end{corollary}
\begin{proof}
For $G\in C(X)$, apply \cref{lem:relative-density-lower-bound} to a solution of ordinary type $\varphi$ with density $c\mathrm{e}^G$, where $c=\mathrm{e}^{B_\varphi(G)}$. It gives $c\|\mathrm{e}^G\|_q\geq b_q$, hence
\[
 B_\varphi(G)\geq\log b_q-\frac1q
                         \log\int_X\mathrm{e}^{qG}\,\mathrm{d}V.
\]

For the upper bound choose a Hermitian form $\Omega\geq\theta$, and let $B_{\Omega,0}$ denote its bounded-type normalization. A bounded solution $v$ for $G$ in this background has a genuine smaller-background envelope. Moreover $u\in\PSH(X,\Omega)$, $u\preceq v$, and positivity gives
\[
 \mathrm{MA}^{\mathrm{ext}}_\Omega(u)
       \geq\mathrm{MA}^{\mathrm{ext}}_\theta(u)
       =c\mathrm{e}^G\,\mathrm{d}V.
\]
If $c>\mathrm{e}^{B_{\Omega,0}(G)}$, the same strict domination argument, now applied to $v$ and $u+t$ in background $\Omega$, is impossible. Hence $B_\varphi\leq B_{\Omega,0}$ and $H_\varphi\leq H_{\Omega,0}$. \Cref{prop:hermitian-dual-entropy-bounds} gives the required upper bound.

For an absolutely continuous probability, take the defining supremum in the lower scalar estimate and use the entropy variational formula with $F=qG$. This proves the lower bound with $a_q=\log b_q$, including infinite entropy. Every non-absolutely-continuous probability has both values infinite by \cref{thm:normalization-dual-compactness}. The finite-domain and bounded-family assertions follow by taking $q=2$. The lower-density lemma is uniform in the reference for fixed $\varepsilon$, and the upper bound uses only $\Omega$. This proves the final assertion.
\end{proof}

For example, let the model references $\varphi_j$ satisfy the same smaller-background condition with fixed $\varepsilon$, let $\|G_j\|_\infty\leq M$, and choose $\nu_j\in\mathcal A_{\varphi_j}(G_j)$. Then
\[
 H_{\varphi_j}(\nu_j)
 =B_{\varphi_j}(G_j)+\int_XG_j\,\mathrm{d}\nu_j
 \leq B_{\Omega,0}(0)+2M.
\]
The uniform lower entropy estimate bounds their entropies. The Hermitian upper estimate places all $\nu_j$ in a single finite sublevel of $H_{\Omega,0}$. Hence \cref{thm:normalization-dual-compactness} gives weak compactness and uniform integrability of their densities. Every weak limit has finite entropy. This gives compactness as the reference varies; identifying the limit as a supporting probability for a limiting reference requires further information.

The corollary retains the model and smaller-background assumptions. It gives no comparison for an arbitrary singular reference outside this solvability regime, and no formula for the extrema of $H_\varphi-\operatorname{Ent}_{\mathrm{d}V}$ when the reference is singular.
